% ============================================================
\chapter{Observable Exhaustivity and the Automorphism Bridge}
\label{ch:dynamics-program}
% ============================================================
\section{Introduction and status}
\label{sec:dyn-intro}
This chapter is a formally marked extension of the main theorem chain.
The preceding chapters establish the closed-system stack through
connection-first geometry, scalar-sector closure, and matter-sector
realization. The present chapter does not supply a premise needed by
those results. It asks the next structural question: how much of the
\emph{state-side} observable content of the closed stack is readable by
the realized geometry, and in what precise sense dynamics can be
recovered from the stack.

The architecture mirrors the interface chapter. There, faithfulness
decomposed into exact failure modes, each defeated by a named
condition, with the realization then canonical. Here, the raw limit
algebra provably exceeds every realized reading
(\cref{thm:dyn-raw-inexhaustible}); the excess is characterized
exactly: it consists of the fiber-separating observables, reached
through two canonical finite-stage mechanisms
(\cref{thm:dyn-ghost-characterization,rem:dyn-ghost-presentations});
a single scale-coherence condition \textup{(U)} selects the
complementary sector, which is provably the maximal sector admitting
a realized reading; and on that sector the realized algebra is
recovered exactly, by construction
(\cref{thm:dyn-regular-identification}). On the same sector,
derivations are classified: they are precisely the vector fields of
the realized manifold (\cref{thm:dyn-derivation-transport}), and the
transport-stabilizing ones form the symmetry algebra of the realized
connection up to gauge (\cref{cor:dyn-stabilizer}).

\begin{remark}[Status ledger]
	\label{rem:dyn-status}
	The following are proved unconditionally relative to the standing
	data \textup{(R1)}--\textup{(R3)} below:
	\cref{lem:dyn-limit-algebra,thm:dyn-raw-inexhaustible,lem:dyn-regular-structure,thm:dyn-regular-identification,lem:dyn-ghost-localization,thm:dyn-ghost-characterization,cor:dyn-smooth-sector,thm:dyn-hadamard,thm:dyn-derivation-transport,lem:dyn-gauge-cocycle,lem:dyn-curvature-class,cor:dyn-stabilizer,cor:dyn-conservation}.
	Condition \textup{(U)} (\cref{def:dyn-scale-coherence}) is a named
	selection: that \emph{some} selection is required is a theorem
	(\cref{thm:dyn-raw-inexhaustible}); that \textup{(U)} selects the
	maximal readable sector is a theorem
	(\cref{thm:dyn-ghost-characterization}); whether the closure
	axioms \emph{force} this selection for the physical sector is open
	(\cref{sec:dyn-open}). Clause \textup{(R2)} is an added state-side
	realization hypothesis introduced by this chapter; its status is
	stated plainly below and its removal is part of open problem
	\textup{(O3)}. The intrinsic characterization of the smooth sector
	and of the stabilization clause is open.
\end{remark}

\paragraph{Standing data.}
Throughout the chapter:
\begin{enumerate}[label=\textup{(R\arabic*)}]
	\item
	      Conditions \textup{(T)} and \textup{(D)} hold
	      (\cref{def:condition-T,def:axiom-D}), and a compatible smooth
	      realization is fixed, faithful on the stabilized carrier; by
	      \cref{thm:canonical-realization,rem:ghost-free-reading} this
	      regime is nonvacuous and the faithfulness is discharged, not
	      assumed.
	\item
	      \emph{State-side realization hypothesis.} Writing
	      \(\Phys_k\) for the stage physical quotients determined by the
	      quotient semantics of \cref{ch:rotation,ch:closure}, with
	      bonding maps induced by the refinement tower and
	      \(\Phys_\infty:=\varprojlim_k\Phys_k\), it is assumed that the
	      realization descends through quotient semantics to a
	      continuous surjection
	      \[
		      \varphi:\Phys_\infty\longrightarrow M
	      \]
	      onto the realized locus \(M\), which is assumed to be a
	      connected smooth manifold of positive dimension, equipped with
	      the realized Riemannian distance \(d\). This clause is an
	      \emph{added interface hypothesis of the present chapter}: it
	      is not asserted to be contained in, or discharged by, the
	      compatibility notion of \cref{ch:interface,ch:connection}.
	      Compactness of \(M\) is automatic, since \(\Phys_\infty\) is
	      compact and \(M=\varphi(\Phys_\infty)\).
	\item
	      A scale assignment \(\varepsilon_k\downarrow0\) is fixed as in
	      \cref{def:scale-completion}.
\end{enumerate}

% ============================================================
\section{The limit comparison algebra}
\label{sec:dyn-algebra}
Each stage is finite, so each stage physical quotient \(\Phys_k\) is a
finite set, and the algebra of admissible stage observables is, by the
quotient-semantics principle that admissible content is exactly
quotient content, the full finite-dimensional algebra of real
functions on \(\Phys_k\):
\[
	\mathsf A_k:=C(\Phys_k).
\]
Bonding maps \(\Phys_{k+1}\to\Phys_k\) induce unital injections
\(\mathsf A_k\hookrightarrow\mathsf A_{k+1}\), and every
\(\mathsf A_k\) pulls back to an algebra of locally constant functions
on \(\Phys_\infty\).

\begin{definition}[Limit comparison algebra]
	\label{def:dyn-limit-algebra}
	The \emph{limit comparison algebra} is the uniform closure
	\[
		\mathsf A
		:=
		\overline{\bigcup_{k}\mathsf A_k}
		\subseteq
		C(\Phys_\infty).
	\]
\end{definition}

\begin{lemma}[The limit algebra is the full continuous algebra]
	\label{lem:dyn-limit-algebra}
	\(\mathsf A=C(\Phys_\infty)\).
\end{lemma}
\begin{proof}
	\(\Phys_\infty\) is an inverse limit of finite discrete spaces,
	hence compact, Hausdorff, totally disconnected, and metrizable
	(the tower is countable). The union \(\bigcup_k\mathsf A_k\) is a
	unital subalgebra of real continuous functions; it separates
	points, since two distinct limit points differ at some finite
	stage and every function on the finite set \(\Phys_k\) is a stage
	observable. The Stone--Weierstrass theorem gives density
	\cite{Rudin1987}.
\end{proof}

\begin{definition}[Realized algebra]
	\label{def:dyn-realized-algebra}
	The \emph{realized algebra} is
	\[
		\mathsf A_{\mathrm{real}}
		:=
		\varphi^*C(M)\subseteq\mathsf A,
	\]
	the pullback of the continuous realized observables. Since
	\(\varphi\) is surjective, \(\varphi^*\) is an isometric unital
	embedding and \(\mathsf A_{\mathrm{real}}\) is closed; in
	particular it already contains the closure of
	\(\varphi^*C^\infty(M)\).
\end{definition}

% ============================================================
\section{Raw inexhaustibility}
\label{sec:dyn-inexhaustible}
The naive exhaustivity statement -- that the limit algebra equals the
realized algebra -- is false, and its failure is structural, not an
artifact of a poor choice of realization.

\begin{theorem}[Raw inexhaustibility]
	\label{thm:dyn-raw-inexhaustible}
	\(\mathsf A_{\mathrm{real}}\subsetneq\mathsf A\). More precisely,
	\(\varphi\) is not injective, and for any two distinct points
	\(p\neq q\) with \(\varphi(p)=\varphi(q)\) there is a finite-stage
	observable separating \(p\) from \(q\) which lies in
	\(\mathsf A\setminus\mathsf A_{\mathrm{real}}\).
\end{theorem}
\begin{proof}
	If \(\varphi\) were injective, it would be a continuous bijection
	from a compact Hausdorff space onto \(M\), hence a homeomorphism;
	but \(\Phys_\infty\) is totally disconnected and \(M\) is connected
	of positive dimension. So \(\varphi\) has a fiber containing two
	distinct points \(p\neq q\). Since \(\Phys_\infty\) is a Stone
	space, there is a clopen set \(C\) with \(p\in C\), \(q\notin C\).
	Every clopen subset of an inverse limit of finite discrete spaces is
	pulled back from a finite stage: a clopen set is compact, hence is
	covered by finitely many cylinder neighborhoods, and these finitely
	many cylinders all refine to one common stage. Therefore the
	indicator \(\mathbf 1_C\) lies in some
	\(\mathsf A_k\subseteq\mathsf A\). Every element of
	\(\mathsf A_{\mathrm{real}}\) is constant on
	\(\varphi\)-fibers, and \(\mathbf 1_C\) is not. Since
	\(\mathsf A_{\mathrm{real}}\) is closed, the inclusion is proper
	even after closure.
\end{proof}

\begin{remark}
	\Cref{thm:dyn-raw-inexhaustible} is the state-side analogue of
	\cref{rem:disconnected-ghost}: it proves that no theorem of the
	closed stack alone can force observable exhaustivity on the raw
	algebra. A selection of the physical sector is therefore mandatory.
	What can be proved is that the excess admits an exact
	characterization, and that one canonical selection -- the maximal
	one admitting a realized reading -- recovers the realized algebra
	identically. That is the content of the next two sections.
\end{remark}

% ============================================================
\section{The regular sector}
\label{sec:dyn-coherence}
\begin{definition}[Realized oscillation; scale coherence]
	\label{def:dyn-oscillation}
	For \(a\in\mathsf A\) and \(\varepsilon>0\), the \emph{realized
		oscillation of \(a\) at scale \(\varepsilon\)} is
	\[
		\operatorname{osc}_\varphi(a;\varepsilon)
		:=
		\sup_{x\in M}\;
		\sup\bigl\{\,
		|a(p)-a(q)|
		:
		p,q\in\varphi^{-1}\bigl(B(x,\varepsilon)\bigr)
		\,\bigr\},
	\]
	where \(B(x,\varepsilon)\) is the realized metric ball. The element
	\(a\) is \emph{scale-coherent} if
	\(\operatorname{osc}_\varphi(a;\varepsilon)\to0\) as
	\(\varepsilon\to0\). The \emph{regular sector} is
	\[
		\mathsf A_U
		:=
		\{\,a\in\mathsf A:\ a\ \text{is scale-coherent}\,\}.
	\]
\end{definition}

\begin{lemma}[Structure of the regular sector]
	\label{lem:dyn-regular-structure}
	\(\mathsf A_U\) is a closed unital subalgebra of \(\mathsf A\)
	containing \(\varphi^*C(M)\).
\end{lemma}
\begin{proof}
	Stability under sums is immediate. For products,
	\[
		|a(p)b(p)-a(q)b(q)|
		\le
		\|a\|\,|b(p)-b(q)|+\|b\|\,|a(p)-a(q)|,
	\]
	so
	\(\operatorname{osc}_\varphi(ab;\varepsilon)\le
	\|a\|\operatorname{osc}_\varphi(b;\varepsilon)
	+\|b\|\operatorname{osc}_\varphi(a;\varepsilon)\).
	For closedness,
	\(\operatorname{osc}_\varphi(a;\varepsilon)
	\le\operatorname{osc}_\varphi(a_j;\varepsilon)+2\|a-a_j\|\),
	so
	\(\limsup_{\varepsilon\to0}\operatorname{osc}_\varphi(a;\varepsilon)
	\le2\|a-a_j\|\to0\). For pullbacks, \(M\) is compact, so
	\(f\in C(M)\) is uniformly continuous and
	\(\operatorname{osc}_\varphi(\varphi^*f;\varepsilon)
	\le\sup_x\operatorname{osc}\bigl(f;B(x,\varepsilon)\bigr)\to0\).
\end{proof}

\begin{definition}[Condition \textup{(U)}: scale-coherent admissibility]
	\label{def:dyn-scale-coherence}
	The closed tower satisfies \emph{condition \textup{(U)}} if the
	physical observable sector is the regular sector \(\mathsf A_U\):
	an admissible limit observable is required to have a definite
	reading in the refinement limit, i.e.\ to be scale-coherent.
\end{definition}

\begin{remark}[Status of condition \textup{(U)}]
	\label{rem:dyn-U-status}
	Three statements with three different standings must be kept
	apart. First, that some selection of the physical sector is
	required is a theorem: \cref{thm:dyn-raw-inexhaustible}. Second,
	that \textup{(U)} is the canonical selection is also a theorem:
	by \cref{thm:dyn-ghost-characterization} below, \(\mathsf A_U\)
	consists exactly of the observables constant on realization
	fibers, hence is the \emph{maximal} sector on which any
	fiber-consistent realized reading can exist; \textup{(U)} selects
	that maximum and nothing less. Third, that the closure axioms
	force this selection for the physical sector is \emph{not} proved
	here; it is open problem \textup{(O1)}. In the vocabulary of the
	factorization criterion for route invariants, an observable
	violating \textup{(U)} is one whose reading depends on the
	refinement route rather than on the realized endpoint; whether
	that observation can be promoted to a derivation of \textup{(U)}
	is exactly \textup{(O1)}. Where the transport side required two
	conditions, \textup{(T)} and \textup{(D)}, the state side requires
	one.
\end{remark}

% ============================================================
\section{The regular-sector identification theorem}
\label{sec:dyn-identification}
The next theorem is the chapter's central result: on the regular
sector, exhaustivity is not an assumption but an identity, with an
explicit inverse to the pullback.

\begin{theorem}[Regular-sector identification]
	\label{thm:dyn-regular-identification}
	\(\mathsf A_U=\varphi^*C(M)\). Moreover the \emph{reading map}
	\[
		\mathcal R:\mathsf A_U\longrightarrow C(M),
		\qquad
		\mathcal R(a)(x):=\lim_{k\to\infty}a(p_{k,x}),
		\quad
		p_{k,x}\in\varphi^{-1}\bigl(B(x,\varepsilon_k)\bigr)
		\ \text{arbitrary},
	\]
	is well defined, independent of all choices, and is an isometric
	unital algebra isomorphism inverse to \(\varphi^*\):
	\[
		\mathcal R\circ\varphi^*=\id_{C(M)},
		\qquad
		\varphi^*\circ\mathcal R=\id_{\mathsf A_U}.
	\]
\end{theorem}
\begin{proof}
	\emph{Existence and independence.}
	Fibers of balls are nonempty since \(\varphi\) is onto. For
	\(j,k\ge K\), both \(p_{j,x}\) and \(p_{k,x}\) lie in
	\(\varphi^{-1}(B(x,\varepsilon_K))\), so
	\[
		|a(p_{j,x})-a(p_{k,x})|
		\le\operatorname{osc}_\varphi(a;\varepsilon_K)\to0.
	\]
	The sequence is Cauchy, and the same estimate shows the limit does
	not depend on the choices.

	\emph{Continuity.}
	For \(x,y\in M\), choose the approximants at stage \(k\); since
	\(B(y,\varepsilon_k)\subseteq B(x,d(x,y)+\varepsilon_k)\), both
	\(p_{k,x}\) and \(p_{k,y}\) lie in
	\(\varphi^{-1}(B(x,d(x,y)+\varepsilon_k))\), whence
	\[
		|\mathcal R(a)(x)-\mathcal R(a)(y)|
		\le
		\limsup_k\operatorname{osc}_\varphi
		\bigl(a;\,d(x,y)+\varepsilon_k\bigr),
	\]
	which tends to \(0\) as \(d(x,y)\to0\) by scale coherence. Thus
	\(\mathcal R(a)\in C(M)\).

	\emph{\(\varphi^*\circ\mathcal R=\id\).}
	Fix \(p\in\Phys_\infty\) and set \(x:=\varphi(p)\). Then \(p\)
	itself lies in \(\varphi^{-1}(B(x,\varepsilon_k))\) for every
	\(k\), so
	\[
		|a(p)-a(p_{k,x})|
		\le\operatorname{osc}_\varphi(a;\varepsilon_k)\to0,
	\]
	giving \(a(p)=\mathcal R(a)(\varphi(p))\). Hence
	\(a=\varphi^*\mathcal R(a)\), and in particular
	\(\mathsf A_U\subseteq\varphi^*C(M)\); with
	\cref{lem:dyn-regular-structure} this proves the set equality.

	\emph{\(\mathcal R\circ\varphi^*=\id\).}
	For \(f\in C(M)\),
	\(\mathcal R(\varphi^*f)(x)=\lim_k f(\varphi(p_{k,x}))\) with
	\(\varphi(p_{k,x})\in B(x,\varepsilon_k)\), and continuity of \(f\)
	gives the limit \(f(x)\).

	\emph{Homomorphism and isometry.}
	Each evaluation \(a\mapsto a(p_{k,x})\) is multiplicative and
	linear, hence so is the limit; \(\mathcal R(1)=1\). Isometry holds
	because \(\varphi^*\) is isometric and the two maps are mutually
	inverse.
\end{proof}

\begin{lemma}[Ghost localization]
	\label{lem:dyn-ghost-localization}
	If \(a\notin\mathsf A_U\), there exist \(\eta>0\), a point
	\(x^*\in M\), and points
	\(p\neq q\in\varphi^{-1}(x^*)\) in a \emph{single fiber}, with
	\[
		|a(p)-a(q)|\ge\eta .
	\]
	Every failure of scale coherence is witnessed, in the limit, by a
	distinction inside one realization fiber.
\end{lemma}
\begin{proof}
	Failure of scale coherence gives \(\eta>0\), scales
	\(\varepsilon_j\downarrow0\), points \(x_j\in M\), and witnesses
	\(p_j,q_j\in\varphi^{-1}\bigl(B(x_j,\varepsilon_j)\bigr)\) with
	\(|a(p_j)-a(q_j)|\ge\eta\). The spaces \(M\) and
	\(\Phys_\infty\) are compact and metrizable
	(\cref{lem:dyn-limit-algebra}), so after passing to a subsequence,
	\(x_j\to x^*\), \(p_j\to p\), and \(q_j\to q\). Continuity of
	\(\varphi\) gives
	\(\varphi(p)=\lim\varphi(p_j)=x^*\), since
	\(\varphi(p_j)\in B(x_j,\varepsilon_j)\) and
	\(\varepsilon_j\to0\); likewise \(\varphi(q)=x^*\). Continuity of
	\(a\) gives \(|a(p)-a(q)|\ge\eta\), and in particular
	\(p\neq q\).
\end{proof}

\begin{theorem}[Ghost characterization]
	\label{thm:dyn-ghost-characterization}
	For \(a\in\mathsf A\) the following are equivalent:
	\begin{enumerate}[label=\textup{(\roman*)}]
		\item \(a\) is scale-coherent: \(a\in\mathsf A_U\);
		\item \(a\) is constant on every fiber of \(\varphi\);
		\item \(a\in\varphi^*C(M)\).
	\end{enumerate}
	Consequently
	\(\mathsf A\setminus\mathsf A_U\) consists exactly of the
	fiber-separating observables, and \(\mathsf A_U\) is the maximal
	subset of \(\mathsf A\) on which a fiber-consistent realized
	reading can exist.
\end{theorem}
\begin{proof}
	\textup{(i)}\(\Rightarrow\)\textup{(ii)}:
	if \(\varphi(p)=\varphi(q)=:x\), then
	\(p,q\in\varphi^{-1}(B(x,\varepsilon))\) for every
	\(\varepsilon>0\), so
	\[
		|a(p)-a(q)|
		\le\operatorname{osc}_\varphi(a;\varepsilon)\to0.
	\]

	\textup{(ii)}\(\Rightarrow\)\textup{(iii)}:
	fiber-constancy defines \(f:M\to\mathbb R\) with
	\(a=f\circ\varphi\). A continuous surjection between compact
	Hausdorff spaces is closed, hence a quotient map; for closed
	\(F\subseteq\mathbb R\),
	\(f^{-1}(F)=\varphi\bigl(a^{-1}(F)\bigr)\) is closed, so \(f\) is
	continuous and \(a=\varphi^*f\).

	\textup{(iii)}\(\Rightarrow\)\textup{(i)} is
	\cref{lem:dyn-regular-structure}.

	For the final statements: the complement description is the
	contrapositive of
	\textup{(i)}\(\Leftrightarrow\)\textup{(ii)}, witnessed in
	localized form by \cref{lem:dyn-ghost-localization}; and any
	assignment of realized values must agree on all points of a fiber,
	so it can be defined at most on the fiber-constant elements, which
	by \textup{(ii)}\(\Leftrightarrow\)\textup{(i)} are exactly
	\(\mathsf A_U\).
\end{proof}

\begin{remark}[One species, two presentations]
	\label{rem:dyn-ghost-presentations}
	At finite stages, fiber-separating observables arise through two
	canonical mechanisms: \emph{direct fiber separators} -- clopen
	indicators splitting two points of a single fiber, as in
	\cref{thm:dyn-raw-inexhaustible} -- and \emph{interface
		presentations} -- clopen indicators \(\mathbf 1_C\) whose realized
	images overlap along a boundary point, where the oscillation is
	persistent. By \cref{lem:dyn-ghost-localization} the second
	mechanism also produces, in the limit, a distinction inside a
	single fiber. The two mechanisms are therefore finite-stage
	presentations of \emph{one} underlying species, the fiber ghost;
	the taxonomy is a characterization
	(\cref{thm:dyn-ghost-characterization}), not a dichotomy.
\end{remark}

\begin{corollary}[Exhaustivity in the selected sector]
	\label{cor:dyn-exhaustivity}
	Under condition \textup{(U)}, the physical observable algebra
	coincides with the realized algebra,
	\(\mathsf A_U=\varphi^*C(M)\), by
	\cref{thm:dyn-regular-identification} -- an identity, not an added
	assumption. The conditionality of the exhaustivity program is
	thereby concentrated in \textup{(U)} itself: that some selection is
	needed is \cref{thm:dyn-raw-inexhaustible}; that this one is
	canonical and maximal is \cref{thm:dyn-ghost-characterization};
	that it is forced is open.
\end{corollary}

\begin{corollary}[Smooth sector]
	\label{cor:dyn-smooth-sector}
	Set \(\mathsf A_U^\infty:=\mathcal R^{-1}\bigl(C^\infty(M)\bigr)
	=\varphi^*C^\infty(M)\).
	Then \(\mathcal R\) restricts to an isomorphism
	\(\mathsf A_U^\infty\cong C^\infty(M)\), and
	\(\mathsf A_U^\infty\) is dense in \(\mathsf A_U\), since
	\(C^\infty(M)\) is uniformly dense in \(C(M)\) by smoothing on the
	compact manifold \(M\).
\end{corollary}

\begin{remark}[Toward an intrinsic smooth sector]
	\label{rem:dyn-smooth-intrinsic}
	\Cref{cor:dyn-smooth-sector} defines the smooth sector through the
	reading. The expected intrinsic characterization -- that
	\(\mathsf A_U^\infty\) consists of the scale-coherent observables
	whose realized difference quotients of every order converge under
	refinement, organized by the transport jet tower of
	\cref{sec:jet-tower} -- is open; see \cref{sec:dyn-open}.
\end{remark}

% ============================================================
\section{Derivations of the regular sector}
\label{sec:dyn-derivations}
\begin{remark}[Why dynamics is densely defined]
	\label{rem:dyn-singer-wermer}
	No nontrivial dynamics can act boundedly on the full algebra: by
	the Singer--Wermer theorem, a bounded derivation of a commutative
	Banach algebra maps into the Jacobson radical, and
	\(C(\Phys_\infty)\) is semisimple, so bounded derivations vanish
	\cite{SingerWermer1955}.
	Dynamics therefore necessarily lives on a dense smooth core; the
	canonical core supplied by the chapter is
	\(\mathsf A_U^\infty\).
\end{remark}

\begin{theorem}[Derivations of the smooth algebra are vector fields]
	\label{thm:dyn-hadamard}
	Let \(D:C^\infty(M)\to C^\infty(M)\) be an \(\mathbb R\)-linear map
	satisfying the Leibniz rule. Then there is a unique smooth vector
	field \(X\) on \(M\) with \(D=\Lie_X\).
\end{theorem}
\begin{proof}
	\emph{Locality.} \(D(1)=D(1\cdot1)=2D(1)\) gives \(D(1)=0\), hence
	\(D\) kills constants. If \(f\) vanishes on a neighborhood of
	\(p\), choose \(\chi\in C^\infty(M)\) with \(\chi\equiv1\) on
	\(\supp f\) and \(\chi(p)=0\); then \(f=\chi f\) and
	\[
		Df(p)=\chi(p)\,Df(p)+f(p)\,D\chi(p)=0 .
	\]
	So \(Df(p)\) depends only on the germ of \(f\) at \(p\).

	\emph{Local form.} In a chart around \(p\) with coordinates
	\(x^1,\dots,x^n\), extended to global smooth functions
	\(\widetilde x^i\) by a bump function, Hadamard's lemma writes any
	smooth \(f\) near \(p\) as
	\[
		f=f(p)+\sum_i(\widetilde x^i-x^i(p))\,g_i,
		\qquad
		g_i\ \text{smooth},\quad g_i(p)=\partial_if(p).
	\]
	For the smooth-manifold form of this local decomposition, see
	\cite{Lee2013}.
	Applying \(D\) at \(p\) and using locality and \(D(1)=0\),
	\[
		Df(p)=\sum_i X^i(p)\,\partial_if(p),
		\qquad
		X^i:=D\widetilde x^i .
	\]
	The functions \(X^i\) are smooth, so \(X:=\sum_iX^i\partial_i\) is
	a smooth vector field on the chart, and the expressions glue
	because they are determined by \(D\). Uniqueness holds because
	coordinate functions separate tangent vectors.
\end{proof}

\begin{theorem}[Derivation transport]
	\label{thm:dyn-derivation-transport}
	The assignments
	\[
		\delta\longmapsto X_\delta,
		\qquad
		\Lie_{X_\delta}=\mathcal R\circ\delta\circ\varphi^*,
	\]
	and
	\[
		X\longmapsto\delta_X:=\varphi^*\circ\Lie_X\circ\mathcal R
	\]
	are mutually inverse Lie-algebra isomorphisms between the
	derivations of \(\mathsf A_U^\infty\) and the smooth vector fields
	on \(M\).
\end{theorem}
\begin{proof}
	If \(\delta\) is a derivation of \(\mathsf A_U^\infty\), then
	\(\mathcal R\circ\delta\circ\varphi^*\) is an \(\mathbb R\)-linear
	Leibniz map of \(C^\infty(M)\), since \(\mathcal R\) is an algebra
	isomorphism with inverse \(\varphi^*\)
	(\cref{thm:dyn-regular-identification,cor:dyn-smooth-sector});
	by \cref{thm:dyn-hadamard} it equals \(\Lie_{X_\delta}\) for a
	unique \(X_\delta\). Conversely \(\delta_X\) is a derivation of
	\(\mathsf A_U^\infty\) by the same conjugation. The two maps are
	mutually inverse because
	\(\mathcal R\circ\varphi^*=\id\) and
	\(\varphi^*\circ\mathcal R=\id\) on the smooth sector. Brackets
	are preserved:
	\[
		[\delta_X,\delta_Y]
		=
		\varphi^*\circ[\Lie_X,\Lie_Y]\circ\mathcal R
		=
		\varphi^*\circ\Lie_{[X,Y]}\circ\mathcal R
		=
		\delta_{[X,Y]} .
	\]
\end{proof}

% ============================================================
\section{The stabilizer and the dynamics bridge}
\label{sec:dyn-bridge}
Let \(\nabla\) denote the realized transport connection of
\cref{ch:christoffel,ch:connection}, with curvature \(R^\nabla\), and
for a vector field \(X\) let \(\Lie_X\nabla\) denote the
\(\End(TM)\)-valued one-form measuring the variation of \(\nabla\)
along the flow of \(X\); in a local frame with connection form
\(\omega\) this is \(\Lie_X\omega\).

\begin{definition}[Symmetry algebra of the realized transport]
	\label{def:dyn-symmetry-algebra}
	\[
		\mathfrak{aut}(M,[\nabla])
		:=
		\bigl\{\,X\in\mathrm{Vect}(M)
		:
		\Lie_X\nabla=\mathrm d^{\nabla}\!A_X
		\quad\text{for some } A_X\in\Gamma(\End(TM))\,\bigr\},
	\]
	the fields moving the realized transport law only by an
	infinitesimal gauge transformation.
\end{definition}

\begin{lemma}[Gauge cocycle; closure under bracket]
	\label{lem:dyn-gauge-cocycle}
	In a local frame, for any \(\End\)-valued form \(B\),
	\[
		\Lie_X(\mathrm d^\nabla B)-\mathrm d^\nabla(\Lie_XB)
		=
		[\Lie_X\omega,\,B] .
	\]
	Consequently, if \(X,Y\in\mathfrak{aut}(M,[\nabla])\) with
	potentials \(A_X,A_Y\), then
	\[
		\Lie_{[X,Y]}\nabla
		=
		\mathrm d^\nabla A_{[X,Y]},
		\qquad
		A_{[X,Y]}
		=
		\Lie_XA_Y-\Lie_YA_X+[A_X,A_Y],
	\]
	so \(\mathfrak{aut}(M,[\nabla])\) is a Lie algebra.
\end{lemma}
\begin{proof}
	With \(\mathrm d^\nabla B=\mathrm dB+[\omega,B]\) and
	\([\Lie_X,\mathrm d]=0\),
	\[
		\Lie_X(\mathrm d^\nabla B)
		=
		\mathrm d(\Lie_XB)+[\Lie_X\omega,B]+[\omega,\Lie_XB]
		=
		\mathrm d^\nabla(\Lie_XB)+[\Lie_X\omega,B].
	\]
	Then, using
	\(\Lie_{[X,Y]}=[\Lie_X,\Lie_Y]\) and
	\(\Lie_X\omega=\mathrm d^\nabla A_X\),
	\(\Lie_Y\omega=\mathrm d^\nabla A_Y\),
	\begin{align*}
		\Lie_{[X,Y]}\omega
		 & =
		\Lie_X(\mathrm d^\nabla A_Y)-\Lie_Y(\mathrm d^\nabla A_X)      \\
		 & =
		\mathrm d^\nabla(\Lie_XA_Y)+[\mathrm d^\nabla A_X,A_Y]
		-\mathrm d^\nabla(\Lie_YA_X)-[\mathrm d^\nabla A_Y,A_X]        \\
		 & =
		\mathrm d^\nabla\bigl(\Lie_XA_Y-\Lie_YA_X+[A_X,A_Y]\bigr),
	\end{align*}
	the last step by the Leibniz rule
	\[
		\mathrm d^\nabla[A_X,A_Y]
		=
		[\mathrm d^\nabla A_X,A_Y]+[A_X,\mathrm d^\nabla A_Y].
	\]
\end{proof}

\begin{lemma}[Curvature class preservation]
	\label{lem:dyn-curvature-class}
	If \(X\in\mathfrak{aut}(M,[\nabla])\) with potential \(A_X\), then
	\[
		\Lie_XR^\nabla=[R^\nabla,A_X].
	\]
	The realized curvature moves inside its gauge class, so the reading
	of every stabilized square class of the carrier
	\(\mathcal K_\infty\) is preserved up to conjugation; under the
	faithful reading
	(\cref{thm:interface-curvature,prop:canonical-universality}) this
	is exactly preservation of the carrier classes.
\end{lemma}
\begin{proof}
	In a frame, \(R^\nabla=\mathrm d\omega+\tfrac12[\omega,\omega]\),
	so
	\[
		\Lie_XR^\nabla
		=
		\mathrm d(\Lie_X\omega)+[\Lie_X\omega,\omega]
		=
		\mathrm d^\nabla(\Lie_X\omega)
		=
		\mathrm d^\nabla\mathrm d^\nabla A_X
		=
		[R^\nabla,A_X],
	\]
	the last identity being the standard
	\(\mathrm d^\nabla\circ\mathrm d^\nabla=\mathrm{ad}_{R^\nabla}\).
\end{proof}

\begin{definition}[Stabilizing derivation]
	\label{def:dyn-admissible-derivation}
	A derivation \(\delta\) of \(\mathsf A_U^\infty\) is
	\emph{stabilizing} if its transported field
	\(X_\delta\) (\cref{thm:dyn-derivation-transport}) lies in
	\(\mathfrak{aut}(M,[\nabla])\). Write
	\(\mathrm{Der}_{\mathrm{stab}}(\mathsf A_U^\infty,\mathcal K_\infty)\)
	for the stabilizing derivations.
\end{definition}

\begin{remark}[Intrinsicality of the stabilization clause]
	\label{rem:dyn-stabilization-intrinsic}
	The clause is stated on the realized side. Its intrinsic content,
	via \cref{lem:dyn-curvature-class}, is that \(\delta\) moves each
	stabilized square class within its conjugation class; under the
	discharged faithful reading the two formulations agree in the
	direction proved above. A purely intrinsic formulation, and the
	converse direction, belong to the open problems
	(\cref{sec:dyn-open}).
\end{remark}

\begin{corollary}[Stabilizer identification]
	\label{cor:dyn-stabilizer}
	Derivation transport restricts to a canonical Lie-algebra
	isomorphism
	\[
		\mathrm{Der}_{\mathrm{stab}}
		\bigl(\mathsf A_U^\infty,\mathcal K_\infty\bigr)
		\;\cong\;
		\mathfrak{aut}(M,[\nabla]).
	\]
\end{corollary}
\begin{proof}
	Immediate from
	\cref{thm:dyn-derivation-transport,def:dyn-admissible-derivation,lem:dyn-gauge-cocycle}.
\end{proof}

\begin{corollary}[Dynamics bridge]
	\label{cor:dyn-bridge}
	In the regime \textup{(T)}+\textup{(D)}+\textup{(U)}, with the
	state-side hypothesis \textup{(R2)}, the stabilizing derivations
	of the selected physical sector are identified with the
	infinitesimal symmetries of the realized transport law up to
	gauge,
	\[
		\mathrm{Der}_{\mathrm{stab}}\cong\mathfrak{aut}(M,[\nabla]),
	\]
	acting on the physical observable algebra through derivation
	transport.
\end{corollary}

\begin{remark}[The conditional upgrade]
	\label{rem:dyn-upgrade}
	If open problems \textup{(O1)} and \textup{(O2)} close -- that is,
	if condition \textup{(U)} is forced from the closure axioms and the
	stabilization clause is formulated and proved intrinsically -- then
	\cref{cor:dyn-bridge} upgrades to the unqualified statement that
	admissible dynamics is reconstructed from the closed stack rather
	than imposed on it. The chapter deliberately does not claim that
	statement yet.
\end{remark}

\begin{corollary}[Conservation of endpoint-determined route invariants]
	\label{cor:dyn-conservation}
	Let \(X\in\mathfrak{aut}(M,[\nabla])\) generate a flow that
	preserves \emph{each} endpoint fiber of a comparison family
	setwise. Then every route invariant satisfying the
	endpoint-factorization criterion is constant along the flow.
\end{corollary}
\begin{proof}
	By the factorization criterion, such an invariant is a function of
	the endpoint data alone, i.e.\ of the endpoint fibers. A flow
	preserving each fiber setwise fixes the arguments of that function,
	hence its value.
\end{proof}

% ============================================================
\section{Open problems}
\label{sec:dyn-open}
\begin{enumerate}[label=\textup{(O\arabic*)}]
	\item
	      \emph{Forcing scale coherence.} Derive condition \textup{(U)}
	      from the closure axioms, plausibly through the factorization
	      criterion read in the scale direction: an observable whose
	      reading depends on the refinement route is not
	      endpoint-determined at scale zero.
	      \Cref{thm:dyn-raw-inexhaustible} shows some selection is
	      mandatory and \cref{thm:dyn-ghost-characterization} shows
	      \textup{(U)} selects the maximal readable sector; the question
	      is whether the physical sector is \emph{forced} to be that
	      maximum, as \textup{(T)} is forced in the Magnus regime
	      (\cref{thm:magnus-torsion-free}).
	\item
	      \emph{Intrinsic smooth sector and stabilization clause.}
	      Characterize \(\mathsf A_U^\infty\) inside \(\mathsf A_U\) by
	      transport-jet data (\cref{sec:jet-tower}), and give the
	      intrinsic formulation, with converse, of the stabilization
	      clause of \cref{def:dyn-admissible-derivation}.
	\item
	      \emph{Spectral reconstruction; removal of \textup{(R2)}.}
	      Upgrade the pair
	      \((\mathsf A_U,\ \text{jet filtration},\ \mathcal K_\infty)\)
	      to a spectral datum recovering \(M\) \emph{with its smooth
	      structure} as a spectrum, in the sense of the commutative
	      reconstruction theorem of noncommutative geometry
	      \cite{Connes2013}. This would
	      eliminate the added state-side hypothesis \textup{(R2)}: the
	      descent map \(\varphi\) would be constructed rather than
	      assumed.
	\item
	      \emph{The Hamiltonian layer.} Pass from the symmetry algebra
	      \(\mathfrak{aut}(M,[\nabla])\) to an evolution law. The
	      natural forcing mechanism is closure of the
	      hypersurface-deformation algebra, which is known to determine
	      the Einstein Hamiltonian on metric phase space
	      (\cref{rem:external-certifications}); the task is to derive
	      the deformation-algebra closure condition itself from the
	      closed stack.
\end{enumerate}

\section{Conclusion}
\label{sec:dyn-conclusion}
The exhaustivity question is now in the same final form as the
interface question before it. The raw limit algebra provably exceeds
every realized reading; the excess is exactly the fiber-separating
sector, reached at finite stages through two presentations of one
ghost species; the scale-coherence condition \textup{(U)} selects the
maximal readable sector; and on that sector the realized algebra, its
smooth core, its derivations, and its transport-stabilizing
derivations are identified by construction:
\[
	\mathsf A_U=\varphi^*C(M),
	\qquad
	\mathrm{Der}\bigl(\mathsf A_U^\infty\bigr)\cong\mathrm{Vect}(M),
	\qquad
	\mathrm{Der}_{\mathrm{stab}}\cong\mathfrak{aut}(M,[\nabla]).
\]
What remains open is concentrated and named: forcing \textup{(U)},
intrinsicizing the smooth sector and the stabilization clause,
spectral reconstruction in place of the state-side hypothesis
\textup{(R2)}, and the Hamiltonian layer.
